\section{Verification purpose}

\textbf{Case ID:} \texttt{\CaseID}. This case verifies the transition from surface fire to passive and active crown fire as wind increases.  The verification question is whether all three wind variants produce the expected crown class and maximum spread rate within their respective tolerances.

\section{Mathematical formulation and numerical implementation}

ELMFIRE compares surface fireline intensity with crown-initiation requirements and then evaluates active-crown spread and critical spread conditions.  The output class encodes 0 for surface-only, 1 for passive crown fire, and 2 for active crown fire.  Wind changes both the surface intensity and crown propagation response.

\section{Derivation of expected behavior}

With fixed canopy and fuel properties, the 1, 5, and 10 mph variants are expected to produce classes 0, 1, and 2 with maximum rates \(1.809\), \(10.06\), and \(35.3~\mathrm{m\,min^{-1}}\), respectively. The sequence \(0<1<2\) and increasing rate are invariants.

\section{Verification metrics and acceptance criteria}

Maximum crown class directly tests the discrete transition state and uses absolute tolerance 0.5 so only the intended integer class passes. Maximum finite \texttt{vs} tests the associated propagation rate with 10 percent relative tolerance. Both metrics must pass for every wind variant.

\section{Simulation configuration and predicted outcome}

All variants use Fuel Model 3, uniform canopy cover 80 percent, canopy height 55, canopy base height 16, canopy bulk density 10 in the GUIDE's encoded raster units, flat terrain, and identical ignition. Only wind speed changes among 1, 5, and 10 mph.

\section{Actual simulation results}

The required evidence is generated only from case-local ELMFIRE outputs. The current postprocessing status is \textbf{\VerificationStatus}; \PassingMetricCount{} of \MetricCount{} required metric rows pass across \VariantCount{} variant(s). A missing output is reported as NOT RUN or NOT EVALUABLE and is not interpreted as verification evidence.


\begin{figure}[H]
  \centering
  \IfFileExists{../figures/whole_domain_result.pdf}{%
    \includegraphics[width=0.82\textwidth,height=0.58\textheight,keepaspectratio]{whole_domain_result.pdf}%
  }{\fbox{\parbox{0.85\textwidth}{No case-local whole-domain ELMFIRE raster is available.}}}
  \caption{Final whole-domain ELMFIRE time of arrival for the displayed variant, in seconds. The figure annotation records the selected case-local GeoTIFF path and coordinate reference system; the generated metrics table, not visual appearance, determines acceptance.}
\end{figure}


\section{Calculated metrics and verification decision}

\begin{center}
\begin{tabular}{@{}llll@{}}
\toprule
Variant / metric & Expected & Calculated & Status \\
\midrule
\MetricRows
\bottomrule
\end{tabular}
\end{center}

The authoritative overall decision is \textbf{\VerificationStatus} and is generated by \path{scripts/postprocess.py} from \path{outputs/metrics.json}; it is not recomputed in this report.

\section{Technical interpretation}

A class failure with correct rate isolates transition logic or canopy-unit handling; coupled class and rate errors indicate wind/crown-fire integration. No scientific conclusion should be drawn when the status is NOT RUN or NOT EVALUABLE.

\section{Scope and limitations}

The test covers one canopy structure and three discrete winds; it does not validate crown-fire correlations against observations. This case provides implementation verification evidence only.
